% !PRINTED 195
\weberpage{195}
%<<CONT
$A, B, C, D$ aus dem Additionstheorem folgert. Wenn wir den Wert des
Multiplikators $n$, zu dem diese Funktionen gehören, durch einen Index
andeuten, so haben wir zunächst
\weq{10}{A_1 = 1, \quad B_1 = 1, \quad C_1 = 1, \quad D_1 = 1;}
ferner aus den Additionsformeln [\S\,44, (16), für $u = v$]:
\weq{11}{
\begin{aligned}
\mathrm{sn}\,2v &= \dfrac{2\,\mathrm{sn}\,v\,\mathrm{cn}\,v\,\mathrm{dn}\,v}{1-\varkappa^2\,\mathrm{sn}^4 v}, \\
\mathrm{cn}\,2v &= \dfrac{\mathrm{cn}^2 v - \mathrm{sn}^2 v\,\mathrm{dn}^2 v}{1-\varkappa^2\,\mathrm{sn}^4 v}, \\
\mathrm{dn}\,2v &= \dfrac{\mathrm{dn}^2 v - \varkappa^2\,\mathrm{sn}^2 v\,\mathrm{cn}^2 v}{1-\varkappa^2\,\mathrm{sn}^4 v}.
\end{aligned}
}
% [source verified, 2026-08-22] The scan distinctly prints coordinate x
% in the x^2 and x^4 factors below, beside the separate modulus \varkappa.
\weq{12}{
\begin{aligned}
A_2 &= 2, \\
B_2 &= 1 - 2x^2 + \varkappa^2x^4, \\
C_2 &= 1 - 2\varkappa^2x^2 + \varkappa^2x^4, \\
D_2 &= 1 - \varkappa^2x^4.
\end{aligned}
}
Wenn wir in (11) $v$ durch $nv$ ersetzen, so folgt:
\weq{13}{
\begin{aligned}
A_{2n} &= 2A_nB_nC_nD_n, \\
B_{2n} &= B_n^2D_n^2 - x^2y^2z^2A_n^2C_n^2, & n \text{ gerade,} \\
&= y^2B_n^2D_n^2 - x^2z^2A_n^2C_n^2, & n \text{ ungerade,} \\
C_{2n} &= C_n^2D_n^2 - \varkappa^2x^2y^2z^2A_n^2B_n^2, & n \text{ gerade,} \\
&= z^2C_n^2D_n^2 - \varkappa^2x^2y^2A_n^2B_n^2, & n \text{ ungerade,} \\
D_{2n} &= D_n^4 - \varkappa^2x^4y^4z^4A_n^4, & n \text{ gerade,} \\
&= D_n^4 - \varkappa^2x^4A_n^4, & n \text{ ungerade,}
\end{aligned}
}
und wenn man in den Additionsformeln des \S\,44 $u = nv$, $v = (n+1)v$ setzt:
\weq{14}{
\begin{aligned}
A_{2n+1} &= y^2z^2A_nD_nB_{n+1}C_{n+1} + A_{n+1}D_{n+1}C_nB_n, & n \text{ gerade,} \\
&= A_nD_nB_{n+1}C_{n+1} + y^2z^2A_{n+1}D_{n+1}C_nB_n, & n \text{ ungerade,} \\
B_{2n+1} &= B_nB_{n+1}D_nD_{n+1} - x^2z^2A_nA_{n+1}C_nC_{n+1}, \\
C_{2n+1} &= C_nC_{n+1}D_nD_{n+1} - \varkappa^2x^2y^2A_nA_{n+1}B_nB_{n+1}, \\
D_{2n+1} &= D_n^2D_{n+1}^2 - \varkappa^2x^4y^2z^2A_n^2A_{n+1}^2.
\end{aligned}
}

Aus diesen Formeln schließt man, \gesperrt{daß die $A$, $B$, $C$, $D$ ganze rationale
Funktionen von $\varkappa^2$ sind, und daß die Zahlenkoeffizienten ganze rationale
Zahlen sind.}

Denn nach (10), (12) hat diese Eigenschaft für $n = 1, 2$ statt und folglich
nach (13), (14) allgemein.
